\chapter{Introduction}
\label{chap:intro}

\begin{chapeau}
Ce chapitre présente le problème médical, l'idée du projet et ce que fait le logiciel, en une page. Le détail vient
dans les chapitres suivants.
\end{chapeau}

\section{Le problème : un os qui pousse de travers}

Chez le chien, il arrive qu'un os de la patte arrière pousse de travers. Le fémur ou le tibia prend alors une forme
courbée ou vrillée : on parle de \emph{déformité angulaire}. Les cas les plus fréquents sont le fémur courbé vers
l'intérieur qui accompagne la luxation de la rotule, le fémur courbé vers l'extérieur des grandes races, le tibia
tordu du teckel (le \emph{pes varus}) et le tibia trop incliné vers l'avant ou vers l'arrière
\cite{brower2017,vezzoni2022,petazzoni2012}. Une déformité peut faire boiter le chien, user ses articulations ou faire
sortir sa rotule de sa place.

Pour corriger, le chirurgien fait une \emph{ostéotomie} : il coupe l'os, remet le morceau du bas dans la bonne
direction, puis le fixe avec une plaque vissée. Tout l'enjeu est de savoir \textbf{où couper, comment couper et de
combien tourner}. Une erreur de quelques degrés peut laisser une déformité, et une torsion mal corrigée peut obliger à
réopérer \cite{worth2019}.

\section{Pourquoi un scanner, et pourquoi des guides imprimés}

La méthode de référence pour préparer une ostéotomie est la méthode \emph{CORA} de Paley \cite{paley2002}. Le
chirurgien trace des lignes sur une radiographie de face et sur une radiographie de profil, puis il en déduit
l'endroit et l'angle de la coupe. Cette méthode a deux limites :
\begin{itemize}
  \item elle se fait à la main, vue par vue, alors qu'une vraie déformité est souvent \emph{oblique}, c'est-à-dire
  visible à la fois de face et de profil \cite{kim2017} ;
  \item une radiographie montre mal la \emph{torsion} de l'os : l'os vrillé sur lui-même et l'os simplement mal posé
  sur la table donnent presque la même image \cite{apelt2005,aper2005}.
\end{itemize}

Le scanner (CT, pour \emph{computed tomography}) donne l'os en trois dimensions. On y voit sa vraie forme et on peut
y mesurer la torsion \cite{dudley2006,alaiyan2021}. L'impression 3D permet en plus de fabriquer des \emph{guides
patient-spécifiques} : des pièces en plastique moulées sur l'os du patient, qui guident la scie et les broches pendant
l'opération. Plusieurs équipes vétérinaires les utilisent déjà, avec une précision de l'ordre de 2\degre{}
\cite{hall2019,jeon2023,bright2024}. Mais la préparation de ces guides se fait encore à la main, dans des logiciels de
conception, par des spécialistes.

\section{Ce que fait CT-CORA3D}

Notre objectif est de faire \textbf{toute la chaîne automatiquement}, du scanner aux pièces prêtes à imprimer, en
montrant et en expliquant chaque étape. Le chirurgien garde la décision : le logiciel propose, mesure et explique ; il
ne décide pas à sa place.

\fig{chaine_du_projet.pdf}{Les grandes étapes du projet, du scanner à la chirurgie. Les étapes 1 à 5 sont faites par
le logiciel ; la dernière est faite à la main.}{fig:chaine-projet}

Concrètement, à partir d'un dossier DICOM, le logiciel :
\begin{enumerate}
  \item reconstruit le volume du scanner et isole chaque os en 3D (segmentation), puis lui donne un nom ;
  \item fabrique un fichier 3D étanche par os (format STL, celui de l'impression 3D) ;
  \item trouve automatiquement les repères anatomiques du fémur ou du tibia et mesure ses angles articulaires et sa
  torsion ;
  \item calcule la déformité en 3D (CORA, angle vrai, plan de la déformité) et la correction ;
  \item choisit une stratégie chirurgicale : coin d'ouverture ou de fermeture, niveau de la coupe, charnière, broches,
  voie d'abord ; il explique chaque choix et le chirurgien peut tout changer ;
  \item construit deux pièces à imprimer : un guide de coupe, qui guide la scie et les broches, et un guide de
  réduction, qui tient les fragments dans la position corrigée prévue ;
  \item vérifie tout par des contrôles indépendants et écrit le mode opératoire et des plans 2D cotés.
\end{enumerate}

Tout cela tient dans une application de bureau à deux onglets : \ui{Segmentation CT} et \ui{Planification CORA}.

\section{Ce que ce projet apporte}

Les briques de base existent séparément (logiciels de segmentation, méthode CORA, guides imprimés). L'apport du projet
est de les \textbf{relier} et d'\textbf{automatiser les décisions} qui sont aujourd'hui faites à la main :
\begin{itemize}
  \item des repères anatomiques trouvés sans aide, avec les définitions des articles, et une correction de
  l'écart entre mesure 2D (radiographie) et mesure 3D (\cref{sec:calibration}) ;
  \item une méthode CORA faite en 3D, qui donne directement le coin oblique unique au lieu de deux corrections
  successives (\cref{chap:planification}) ;
  \item une stratégie chirurgicale choisie par des règles écrites et testées sur des centaines de cas
  (\cref{sec:strategie}) ;
  \item un guide conçu par modélisation implicite, qui respecte des contraintes de pose (pas de contre-dépouille, pas
  de matière entre deux os), accompagné d'un guide de réduction (\cref{chap:guides}) ;
  \item une validation systématique par des déformités simulées dont on connaît la réponse
  (\cref{chap:validation}).
\end{itemize}

\section{Le cadre}

\begin{itemize}
  \item \textbf{Projet} : PMD2, master en ingénierie biomédicale à l'ECAM Bruxelles, 2026--2027.
  \item \textbf{Données} : trois scanners d'un lévrier sain (hanche--genou, genou--cheville, pied). Les déformités
  sont donc simulées sur des os sains. Un premier cas clinique a été reçu fin octobre 2026 (\cref{sec:topogramme}).
  \item \textbf{Membre traité} : le membre pelvien (fémur et tibia). Au fémur, seulement les déformités proches du
  genou (\cref{sec:hypotheses}).
  \item \textbf{Statut} : prototype de recherche. Il ne doit pas servir sur un animal sans validation par le
  chirurgien.
\end{itemize}

\section{Plan du document}

La \textbf{partie I} pose le problème : les notions d'anatomie, de chirurgie et d'imagerie (\cref{chap:notions}),
puis l'état de l'art (\cref{chap:etat}). La \textbf{partie II} suit le chemin d'un cas, du scanner au guide : vue
d'ensemble (\cref{chap:chaine}), lecture du DICOM (\cref{chap:lecture}), segmentation (\cref{chap:segmentation}),
identification des os (\cref{chap:identification}), maillage et contrôle (\cref{chap:maillage}), repères
(\cref{chap:reperes}), planification CORA (\cref{chap:planification}) et guides (\cref{chap:guides}). Chaque chapitre
donne le principe, les calculs, puis le fichier et la fonction du code qui les font. La \textbf{partie III} décrit
l'application (\cref{chap:application}), la validation (\cref{chap:validation}), les limites et la suite
(\cref{chap:limites}). Les annexes détaillent le code fichier par fichier, tous les paramètres, le format de la
prescription, un glossaire et les commandes utiles.
